\begin{abstract}
GPU resource management policies govern how a system places memory and schedules GPU tasks.
% GPU 资源管理策略决定系统如何放置内存并调度 GPU 任务。
They play an important role in determining application performance, and the ideal policy for a given deployment often depends on low-level details such as the application, workloads, and hardware.
% 它们对应用性能有重要影响，而给定部署场景下的理想策略往往取决于应用、负载和硬件等底层细节。
AI agents would make it possible to automatically optimize GPU resource management policies, but require three properties of a system: safety, to ensure that a broken policy cannot cause a system crash; dynamism, to enable policy changes without restarts; and full-stack, to ensure that the policy can observe device and application states and control low-level hardware.
% AI 智能体有望自动优化 GPU 资源管理策略，但系统需具备三个特性：安全性，确保有缺陷的策略不会导致系统崩溃；动态性，支持策略在不重启的情况下变更；全栈性，确保策略能观测设备与应用状态并控制底层硬件。
Unfortunately, existing techniques for GPU resource management fail to meet these needs.
% 遗憾的是，现有的 GPU 资源管理技术无法满足这些需求。

We introduce \sys, a new system that provides an OS interface for GPU resource management that is safe, dynamic, and full-stack by turning to eBPF: \sys{} policies include eBPF callback functions that the system invokes on a specified event (e.g., when code reaches a function).
% 我们提出 gpubpf，一个借助 eBPF 提供安全、动态、全栈的 GPU 资源管理 OS 接口的新系统：gpubpf 策略包含 eBPF 回调函数，系统在指定事件（如代码执行到某个函数）时调用这些函数。
The system includes a novel execution model that separates the effect of a callback from its triggering event.
% 该系统包含一种新颖的执行模型，将回调的效果与其触发事件分离。
Namely, \sys{} supports policies that asynchronously transition resources between states, potentially invoking multiple transitions at a single point-in-time.
% 具体而言，gpubpf 支持策略在资源状态之间进行异步迁移，并可能在同一时间点触发多次迁移。
We implemented \sys{} on the Linux NVIDIA GPU driver and find that AI-generated policies improve throughput by $1.76\times$ over application-tuned baselines and reduce tail latency by up to $2\times$, without application modifications. The \sys{} mechanism incurs 0.7--3.3\% overhead when executing policies reproduced from prior work.
% 我们在 Linux NVIDIA GPU 驱动上实现了 gpubpf，发现 AI 生成的策略在不修改应用的前提下，吞吐量较应用调优的基线提升 $1.76\times$，尾延迟最多降低 $2\times$。
Moreover, we find that there were zero OS kernel panics through the optimization process, highlighting \sys{}'s safety.
% 此外，整个优化过程中 OS 内核零崩溃，凸显了 gpubpf 的安全性。
\end{abstract}
